\section{Current Operating-Model Evidence}

Team 4 supplies the inherited operating scenarios, refreshed with Q1--Q2 2026 statistics. Its earlier gold simulation is excluded: all price paths come from Team 8. These fixed inputs permit model attribution but are not independently validated operating forecasts.

\subsection{Cost-of-sales block}
The cost model in Section~3.4 extracts a common log-space trend, forecasts it with ARIMA(3,1,0) with drift, and reconstructs mine-level paths with volatility-scaled anchors. Its reported selection and hold-out diagnostics are inherited evidence, not revalidated forecasts or a jointly simulated cost state.

\noindent\begin{minipage}{\columnwidth}
{\centering
  \includegraphics[width=0.95\columnwidth]{figures/team4_mine_cost_qoq.png}
  \captionof{figure}{Selected-mine Cost of Sales per ounce in Q1 and Q2 2026. Mine dispersion motivates preserving the Team 4 operating decomposition rather than replacing it with a single company-average cost.}
  \label{fig:team4_cost_qoq}
\par}
\end{minipage}\par

The scenario uses 20 fixed quarterly cost entries, including two 2026 actuals and 18 forecast entries. It does not propagate the source study's operating distributions.

\subsection{Production block}
Production is reconstructed from ore processed, grade and recovery using the specifications in Section~3.4. This preserves mine heterogeneity. The SARIMA orders are admitted from the source study, which does not report their selection by AIC/BIC.

\noindent\begin{minipage}{\columnwidth}
{\centering
  \includegraphics[width=0.95\columnwidth]{figures/team4_process_components_qoq.png}
  \captionof{figure}{Ore processed, grade and recovery updated with Barrick 2026 actuals. These are physical operating drivers; no commodity-price simulation enters the panel.}
  \label{fig:team4_process_qoq}
\par}
\end{minipage}\par

\noindent\begin{minipage}{\columnwidth}
{\centering
  \includegraphics[width=0.95\columnwidth]{figures/team4_mine_production_qoq.png}
  \captionof{figure}{Selected-mine attributable gold production in Q1 and Q2 2026. The current actuals reveal heterogeneous changes that are hidden by a company-level total.}
  \label{fig:team4_production_qoq}
\par}
\end{minipage}\par

The original study also develops Goodman variance propagation and confidence bands for the production components. Those distributions remain useful evidence on operating-model uncertainty, but they are not propagated jointly in the four-engine corporate experiment: only the refreshed point vectors are passed downstream.

\subsection{Forecast handoff used by the valuation}
The final handoff keeps observed and modelled quantities distinct. Barrick's Q1--Q2 2026 actual production and CoS establish the current operating state; the Team 4 forecasts begin at Q3 2026 and extend through the common valuation horizon.

\noindent\begin{minipage}{\columnwidth}
{\centering
  \includegraphics[width=0.95\columnwidth]{figures/team4_operating_contract.png}
  \captionof{figure}{Inherited operating contract. Company-wide Q1--Q2 2026 actuals precede the eight-mine Team 4 forecast from Q3; the scope break is disclosed in Section 2.5. The straight cost path is the conditional mean of a persistent-trend specification, not a claim that operating uncertainty disappears.}
  \label{fig:team4_contract}
\par}
\end{minipage}\par

\begin{methodbox}[Prudential operating boundary]
Option-implied laws supply conditional scenario shapes to the proxy layer. Costs retain a persistent trend without additional technology-driven cost reductions. This is prudential relative to otherwise identical efficiency-improvement scenarios, but it does not establish a mathematical lower bound. Stochastic costs could yield favourable and adverse outcomes; the current fixed vectors preserve attribution across gold engines. Growth is already included in the common DCF schedule.
\end{methodbox}
